\documentclass[10pt,a4paper]{amsart}
\usepackage[margin=19mm]{geometry}
\usepackage{amsmath,amssymb,mathtools}
\usepackage[scr=rsfs,cal=euler]{mathalfa}
\usepackage{xfrac}
\usepackage[unicode,hidelinks,bookmarks=false]{hyperref}
\newcommand{\ie}{i.e.\ }
\newcommand{\cf}{cf.\ }
\newcommand{\resp}{respectively\ }
\newcommand{\Subsection}{\subsection}
\providecommand{\nobreakdash}{\nobreak\hbox{-}}
\setlength{\emergencystretch}{2em}
\multlinegap0pt
\usepackage{amssymb}
\usepackage[shortlabels]{enumitem}
\usepackage{float}
\usepackage{mathtools}
\usepackage{tikz}
\usepackage{xcolor}
\newtheorem{proposition}{Proposition}
\newcommand{\SR}{S}
\renewcommand{\ge}{\geqslant}
\renewcommand{\le}{\leqslant}
\newcommand{\nf}{\infty}
\title{Condition numbers of block Toeplitz matrices and stability of space-time IgA approximations for the wave and Schrödinger equations}
\author{Manuel Bogoya, Albrecht B\"ottcher, Matteo Ferrari, Sergei M Grudsky, Stefano Serra-Capizzano}
\date{Source: arXiv:2608.24151v1 (CC BY 4.0)}
\begin{document}
\maketitle

\noindent\emph{Source and licence.} Condition numbers of block Toeplitz matrices and stability of space-time IgA approximations for the wave and Schrödinger equations. Version arXiv:2608.24151v1 (CC BY 4.0), distributed by arXiv under the Creative Commons Attribution 4.0 International licence, http://creativecommons.org/licenses/by/4.0/, as recorded in the arXiv licence metadata for this exact version and shown on the abstract page https://arxiv.org/abs/2608.24151v1. Full licence notice: \texttt{licenses/arxiv-2608.24151-CC-BY-4.0.txt}.\par\smallskip

\noindent\textit{Editorial scope.} Counted unit: the article's proposition prop:43 (source0.tex lines 1702--1708). The article's complete original proof of it is delivered with it (source0.tex lines 1710--1764, one \texttt{proof} environment of 55 lines). No mathematics is written, repaired or re-derived: the statement and the proof are the article's own lines, in the article's own notation, and no retained line is substituted or re-flowed.  Retained uncounted as the source-local context the proof needs: lines 1652--1695.  Nothing was omitted from any retained span, so the package carries no omission record.  No retained line contains a citation, so the wrapper carries no bibliography and no bibliography entry.  No retained line contains a figure environment, an image-inclusion command or drawing code, so no image is delivered and the package declares no asset dependency.  Editorial formatting only, in addition: an AMS \texttt{amsart} preamble that restates the article's own declarations and macros used by the retained lines, namely the environments \texttt{proposition} on separate counters, together with the standard packages those lines need.  The retained environments print in this document as Proposition 1 in order of appearance, whereas the article prints the same items with the numbering of its own full document.  The complete native source of the article as distributed in the arXiv e-print for this exact version is bundled unchanged as \texttt{sources/208452/source0.tex}.
\begin{align}\label{eq:422}
\widetilde{\mathbf{M}}_{n}^{(3,1)} &\to
\overbrace{\frac{1}{560}\begin{pmatrix}[r]
0 & 9
\\ 0 & 1
\end{pmatrix}}^{2}
\quad
\overbrace{\frac{1}{560}\begin{pmatrix}[r]
53 & 128
\\ 9 & 80
\end{pmatrix}}^{1}
\quad
\overbrace{\frac{1}{560}\begin{pmatrix}[r]
80 & 9
\\ 128 & 53
\end{pmatrix}}^{0}
\quad
\overbrace{\frac{1}{560}\begin{pmatrix}[r]
1 & 0
\\ 9 & 0
\end{pmatrix}}^{-1}
\\
\widetilde{\mathbf{B}}_{n}^{(3,1)} &\to
\overbrace{\frac{3}{40}\begin{pmatrix}[r]
0 &-5
\\ 0 &-1
\end{pmatrix}}^{2}
\quad
\overbrace{\frac{3}{40}\begin{pmatrix}[r]
-5 & 16
\\-5 & 0
\end{pmatrix}}^{1}
\quad
\overbrace{\frac{3}{40}\begin{pmatrix}[r]
0 &-5
\\ 16 &-5
\end{pmatrix}}^{0}
\quad
\overbrace{\frac{3}{40}\begin{pmatrix}[r]
-1 & 0
\\-5 & 0
\end{pmatrix}}^{-1}
\label{eq:423}
\end{align}

\begin{proposition}\label{prop:43}
The polynomial $\SR_{\rho}^{(3,1)}$ is of type
\begin{align*}
&(3,0,2)\quad\text{if and only if}\quad\rho\in (168/17,10)\cup (42,+\nf),\\
&(2,2,1)\quad\text{if and only if}\quad\rho\in [0,168/17]\cup [10,42].
\end{align*}
\end{proposition}

\begin{proof}
We compute explicitly, using~\eqref{eq:422} and~\eqref{eq:423},
\[
\det \SR_{\rho}^{(3,1)}(t)=\frac{1}{11200}\, t \big(a_{\rho} t^{4}+b_{\rho} t^{3}+c_{\rho} t^{2}+b_{\rho} t+a_{\rho}\big).
\]
Here,
\[
a_{\rho}=-\rho^{2}-48\rho-1260,\quad
b_{\rho}=72\rho^{2}-768\rho+10080,\quad
c_{\rho}=-262\rho^{2}+6672\rho-17640.
\]
We analyze the roots of the reciprocal polynomial
\[
p_{\rho}(t)=a_{\rho} t^{4}+b_{\rho} t^{3}+c_{\rho} t^{2}+b_{\rho} t+a_{\rho}.
\]
We apply the substitution $y=t+t^{-1}$, which yields $t^{2}+t^{-2}=y^{2}-2$.
This reduces the quartic equation to the quadratic equation
\[
f_{\rho}(y)=a_{\rho} y^{2}+b_{\rho} y+(c_{\rho}-2a_{\rho})=0.
\]
Observe that if $y$ is real and $|y|>2$, then the equation $t^{2}-y t+1=0$ has two real roots, one strictly inside and one strictly outside the unit circle.

Conversely, if $y$ is real and $|y|\le 2$, the roots $t$ are complex conjugates lying on the unit circle.
Hence, there is a direct correspondence between pairs of roots of $\det \SR_{\rho}^{(3,1)}$ on the unit circle and roots of $f_{\rho}$ in the interval $[-2,2]$.

We now analyze $f_{\rho}(y)$.
First, note that the leading coefficient satisfies $a_{\rho}<0$ for all $\rho>0$. Moreover, $f_{\rho}(4)=12(\rho^{2}+244\rho+420)>0$ for all $\rho>0$, so that, since $a_{\rho}<0$, the parabola $f_{\rho}$ has two distinct real roots $y_{1}\le y_{2}$, with $y_{2}>4>2$.
Next, we show that $-b_{\rho}/(2a_{\rho})>2$ for all $\rho>0$, which is equivalent to
\[
4a_{\rho}+b_{\rho}=68\rho^{2}-960\rho+5040>0.
\]
The discriminant of this quadratic is $-449280<0$, so the inequality holds for all $\rho>0$.
Therefore, the critical point $-b_{\rho}/(2a_{\rho})$ of $f_{\rho}$ lies strictly to the right of $y=2$.

Since the parabola opens downward and its maximum is attained at $-b_{\rho}/(2a_{\rho})>2$, both roots $y_{1},y_{2}$ are strictly greater than $2$ if and only if $f_{\rho}(2)<0$.
We compute
\[
f_{\rho}(2)=-120\rho(\rho-42).
\]
Thus, for $\rho\in (42,+\nf)$, the polynomial $f_{\rho}$ has two real roots greater than $2$, and hence $\SR_{\rho}^{(3,1)}$ is of type $(3,0,2)$.

We now consider the case $\rho\le 42$.
In this range, $f_{\rho}(2)\ge 0$, so that $y_{1}\le 2\le y_{2}$.

To locate the root $y_{1}$, we evaluate at $y=-2$:
\begin{align*}
f_{\rho}(-2)=-24(\rho-10)(17\rho-168).
\end{align*}
For $\rho\in [0,168/17]\cup [10,42]$, we have $f_{\rho}(-2)\le 0$, and hence $y_{1}\in [-2,2]$.
Therefore, $\SR_{\rho}^{(3,1)}$ is of type $(2,2,1)$.

For $\rho\in (168/17,10)$, we have $f_{\rho}(-2)>0$.
Since also $f_{\rho}(2)>0$ and the parabola opens downward, both roots lie outside the interval $[-2,2]$.
Thus, $\SR_{\rho}^{(3,1)}$ is of type $(3,0,2)$.
\end{proof}

\end{document}
